\documentclass[12pt]{article}

\usepackage{sbc-template}
\usepackage{graphicx,url}
\usepackage[T1]{fontenc}
\usepackage[utf8]{inputenc}
\usepackage[brazil]{babel}

\sloppy

\title{Mineração de Métricas DORA em Repositórios Open-Source que Utilizam GitHub Actions}

\author{Pedro Maia Alves\inst{1}, Diogo Brunoro\inst{1},
  Lorran [PREENCHER SOBRENOME]\inst{1}}

\address{Engenharia de Software -- Pontifícia Universidade Católica de Minas Gerais
  (PUC Minas)\\
  \email{PREENCHER, PREENCHER, PREENCHER}}

\begin{document}

\maketitle

\begin{abstract}
  This study investigates delivery performance in popular open-source repositories
  that use GitHub Actions. Releases, commits, and workflow runs are adopted as
  explicit operational proxies for the four DORA metrics. The study also evaluates
  how alternative definitions affect repository classification. This document
  presents the research questions and hypotheses established before data analysis.
\end{abstract}

\begin{resumo}
  Este estudo investiga o desempenho de entrega em repositórios open-source populares
  que utilizam GitHub Actions. Releases, commits e execuções de workflows são adotados
  como proxies operacionais explícitos para as quatro métricas DORA. O estudo também
  avalia como definições alternativas afetam a classificação dos repositórios. Este
  documento apresenta as questões de pesquisa e as hipóteses definidas antes da análise.
\end{resumo}

\section{Introdução}
\label{sec:introducao}

As métricas DORA (\textit{DevOps Research and Assessment}) tornaram-se o padrão de mercado para
medir o desempenho de entrega de software~\cite{forsgren2018accelerate,dora}. Quatro métricas
descrevem duas dimensões: velocidade, com a frequência de implantação (\textit{deployment
frequency}) e o tempo de entrega de mudanças (\textit{lead time for changes}); e estabilidade,
com a taxa de falha das mudanças (\textit{change failure rate}, CFR) e o tempo de recuperação.

Medir essas métricas em projetos \textit{open-source} esbarra em um problema de construto: o GitHub
não registra implantações em produção nem falhas em produção. O que existe são \textit{releases},
\textit{commits} e execuções de \textit{workflows} de integração contínua. Neste trabalho, esses
artefatos funcionam como \textit{proxies}: uma \textit{release} publicada representa uma
implantação, e uma execução de CI que falhou representa uma falha. Um \textit{proxy} pode enganar.
Uma biblioteca que publica uma \textit{release} por mês não implanta nada em produção, pois quem
faz isso são os seus usuários, e uma falha de pipeline não é uma falha em produção.

O objetivo deste estudo é minerar automaticamente as métricas DORA em repositórios populares do
GitHub que usam GitHub Actions e avaliar quanto se pode confiar nessas medidas. Para isso, o
grupo (i) calcula as métricas com definições operacionais explícitas, (ii) valida os critérios
automáticos contra o julgamento humano e (iii) mede o quanto as conclusões mudam quando a
definição da métrica muda. Este artigo responde às seguintes questões de pesquisa (RQ):

\begin{itemize}
  \item \textbf{RQ01:} qual a frequência de \textit{deploys} dos repositórios populares que usam CI/CD?
  \item \textbf{RQ02:} qual o tempo entre um \textit{commit} e o seu respectivo \textit{deploy}?
  \item \textbf{RQ03:} qual a taxa de falha das mudanças entregues por esses repositórios?
  \item \textbf{RQ04:} qual o tempo de recuperação após uma execução de CI/CD com falha?
  \item \textbf{RQ05:} repositórios com maior frequência de \textit{deploy} apresentam maior ou menor taxa de falha?
  \item \textbf{RQ06:} quais características dos repositórios estão associadas a um melhor desempenho DORA?
  \item \textbf{RQ07:} o quanto a classificação DORA de um repositório depende da definição operacional escolhida?
\end{itemize}

\subsection{Hipóteses}
\label{sec:hipoteses}

As hipóteses abaixo foram escritas antes de qualquer análise dos dados coletados.

\begin{itemize}
  \item \textbf{H1 (RQ01).} Esperamos que a mediana fique abaixo de uma \textit{release} por
  semana, pois a maioria dos projetos populares são bibliotecas e ferramentas que publicam versões
  por ciclos de funcionalidades, e não de forma contínua. Esperamos também uma cauda longa, com
  poucos repositórios publicando \textit{releases} diariamente.
  \item \textbf{H2 (RQ02).} Esperamos que o \textit{lead time} por \textit{release} (variante a)
  seja bem maior que o por \textit{commit} (variante b), porque a variante (a) é dominada pelo
  \textit{commit} mais antigo de cada \textit{release}, enquanto a (b) pondera todos os
  \textit{commits}, a maioria deles escrita perto da publicação.
  \item \textbf{H3 (RQ03).} Esperamos que o CFR de pipeline (variante a) seja maior que o CFR de
  entrega (variante b), pois falhas de CI são baratas, frequentes e ocorrem justamente para
  impedir que o problema chegue a uma \textit{release}, enquanto \textit{releases} corretivas em
  até sete dias são mais raras.
  \item \textbf{H4 (RQ04).} Esperamos que a mediana do tempo de recuperação seja de poucas horas,
  já que falhas de CI costumam ser corrigidas por um novo \textit{commit} no mesmo dia, mas com
  cauda pesada e uma proporção não desprezível de episódios censurados, de \textit{workflows}
  quebrados que ficam abandonados por semanas.
  \item \textbf{H5 (RQ05).} Esperamos uma correlação fraca e negativa entre a frequência de
  \textit{deploys} e o CFR, isto é, que repositórios que entregam com maior frequência tenham
  taxa de falha igual ou menor. Essa expectativa segue a evidência do DORA de que velocidade e
  estabilidade não constituem necessariamente um \textit{trade-off}.
  \item \textbf{H6 (RQ06).} Esperamos que o tipo do projeto e o tamanho da comunidade de
  contribuidores apresentem associações mais claras com o desempenho DORA do que a linguagem
  principal. Aplicações e serviços, por automatizarem entregas aos usuários, devem publicar com
  maior frequência, enquanto comunidades maiores devem se recuperar mais rapidamente de falhas.
  \item \textbf{H7 (RQ07).} Esperamos concordância moderada ou alta entre as classificações
  produzidas pelas diferentes definições operacionais. As mudanças de categoria devem se
  concentrar nos repositórios próximos aos limites das faixas, especialmente quando tags e
  pré-releases forem usadas como unidades de entrega.
\end{itemize}


\bibliographystyle{sbc}
\bibliography{referencias}

\end{document}
